\label{anlage:katalog}

Diese Anlage dokumentiert die im Berichtszeitraum entwickelten Oberflächen- und Navigationsstände des Bauteilkatalogs (\hyperref[abschnitt:human-interface-design]{Abschnitt~\ref*{abschnitt:human-interface-design}}). Die Screenshots zeigen Zwischenstände der Benutzeroberfläche und keinen abgeschlossenen Funktionsumfang.

\begin{center}
\textbf{Bauteilobjekt}\\[2mm]
\includegraphics[width=5.5cm]{asset/katalog/🖼️bauteilobjekt.png}
\end{center}

Die Katalognavigation führt von der Typologienübersicht über die Typen einer Struktur bis zu den einzelnen Bauteilen (Pieces). Das aktive Bauteilobjekt ABAUF-CB-001 zeigt Anschlüsse, den Nachweisstatus je Fachbereich (Geometrie, Statik, Brandschutz, System, CO\textsubscript{2}) sowie den Reifegrad.

\begin{center}
\textbf{Filteransicht}\\[2mm]
\includegraphics[width=13.5cm]{asset/katalog/🖼️filteransicht.png}
\end{center}

Das Filter-Panel gliedert Filterkriterien nach Typologie, Materialfamilie, Geometrie, Funktion im Entwurf, Verfügbarkeit, Reifegrad sowie Risiko und Nachweise. Aktive Filter werden als Chips oberhalb der Kriteriengruppen zusammengefasst.

\begin{center}
\textbf{Bauteilkarte}\\[2mm]
\includegraphics[width=13.5cm]{asset/katalog/🖼️bauteilkarte.png}
\end{center}

Der hierarchische Katalogbaum gliedert Bauteile nach Systemebene (Structure, Envelope, Interior) und darunterliegenden Typen; die Bauteilkarte am unteren Rand bündelt Geometrie, Anschlüsse und Nachweise des ausgewählten Bauteils.
